% Generated draft: requires mathematical and source review.
\documentclass[11pt]{article}
\usepackage[margin=25mm]{geometry}
\usepackage{fontspec}
\setmainfont{Latin Modern Roman}
\usepackage{amsmath,amssymb,mathtools,unicode-math}
\setmathfont{Latin Modern Math}
\usepackage{enumitem,needspace,xurl,hyperref}
\setlength{\parindent}{0pt}
\setlength{\parskip}{5pt}
\setlength{\emergencystretch}{3em}
\newcommand{\N}{\mathbb{N}}
\newcommand{\Z}{\mathbb{Z}}
\newcommand{\Q}{\mathbb{Q}}
\newcommand{\R}{\mathbb{R}}
\newcommand{\C}{\mathbb{C}}
\newcommand{\F}{\mathbb{F}}
\newcommand{\A}{\mathbb{A}}
\newcommand{\PP}{\mathbb{P}}
\newcommand{\E}{\mathbb{E}}
\newcommand{\eps}{\varepsilon}
\newcommand{\dd}{\,\mathrm{d}}
\begin{document}
\section*{UnsolvedMath ID: 800420}
\textbf{Source problem code:} catalogue entry
\subsection{Definitions and mathematical statement}
Let \(q\) be a prime power and let \(\F_q\) denote the field with \(q\) elements. For a polynomial or rational function \(f\), write \(f^{\circ n}\) for its \(n\)-fold compositional iterate, with \(f^{\circ 1}=f\). Thus the source notation \(f^n\) denotes iteration, not an ordinary power.

In the polynomial setting, the source takes \(f\in\F_q[x]\) to be an \emph{(irreducible)} polynomial and takes \(g\in\F_q[x]\). The parentheses around ``irreducible'' leave unclear whether irreducibility is required or is a special case of interest. Retain this qualification; no assumption \(\deg f\ge2\) is supplied.

For each integer \(n\ge1\) for which \(h_n=g\circ f^{\circ n}\) is nonzero and nonconstant, write its normalized factorization as
\[
 h_n(x)=c_n\prod_{i=1}^{r_n}p_{n,i}(x)^{e_{n,i}},
 \qquad c_n\in\F_q^\times,
\]
where the \(p_{n,i}\) are pairwise distinct monic irreducible polynomials and the \(e_{n,i}\) are positive integers. The scalar \(c_n\) accounts for the unit omitted in the displayed factorization in the source and does not affect the quantities below. Define
\[
 R_n=r_n,\qquad
 M_n=\sum_{i=1}^{r_n}e_{n,i},\qquad
 E_n=\max_{1\le i\le r_n}e_{n,i},\qquad
 D_n=\max_{1\le i\le r_n}\deg p_{n,i}.
\]
The source asks the following questions for the given data \(q,f,g\):
\begin{enumerate}
\item[(a)] Describe the behavior of quantities such as \(R_n,M_n,E_n,D_n\) as \(n\to\infty\).
\item[(b)] Given a factorization of \(g\circ f^{\circ n}\), describe or predict the factorization of \(g\circ f^{\circ(n+1)}\).
\item[(c)] Formulate and answer the analogous questions when \(f\in\F_q(x)\) is a rational function. For each defined, nonzero rational function \(h_n=g\circ f^{\circ n}\), use its reduced factorization
\[
 h_n(x)=c_n\prod_{i=1}^{r_n}p_{n,i}(x)^{e_{n,i}},
 \qquad c_n\in\F_q^\times,\quad
 e_{n,i}\in\Z\setminus\{0\},
\]
with pairwise distinct monic irreducible \(p_{n,i}\in\F_q[x]\). Positive exponents record numerator factors and negative exponents record denominator factors. The source suggests that it may be more natural to study
\[
 A_n=\sum_{i=1}^{r_n}|e_{n,i}|,
 \qquad
 B_n=\max_{1\le i\le r_n}|e_{n,i}|
\]
in place of the signed sum and the maximum signed exponent. The analogous questions also concern \(r_n\), the largest factor degree, and prediction of the next factorization.
\end{enumerate}

\textit{Unresolved scope and boundary conditions.} These are requests for a description or prediction, rather than a specified conjectural formula, bound, or limiting law. The source does not exclude \(g=0\) or constant \(g\). Zero has no factorization of the displayed kind; a nonzero constant has an empty factorization, for which the count and empty sum are zero but the displayed maxima have no supplied convention. The definitions involving maxima above therefore apply only when \(r_n\ge1\). In the rational setting, the source does not explicitly state whether \(g\) remains polynomial or may also be rational, nor does it specify how to handle constant \(f\), undefined compositions, or zero compositions. The rational factorization above describes finite irreducible polynomial factors; the source does not separately specify whether the point at infinity should be included in an alternative projective formulation. No additional hypotheses or conventions resolving these issues are asserted here.
\subsection{Short English statement}
For a polynomial over a finite field, study how the number, multiplicities, and degrees of the irreducible factors of \(g\circ f^{\circ n}\) change under iteration, and how one factorization predicts the next. Ask analogous questions for rational functions, allowing negative exponents for denominator factors and considering absolute multiplicities. The intended scope of irreducibility, constant inputs, the rational-function domain of \(g\), and the required form of a description or prediction remain unspecified.
\subsection{Sources}
\begin{itemize}
\item Local UnsolvedMath registry: permanent ID \(800420\), exact source problem code \texttt{catalogue entry}, title ``Open Problems in Arithmetic Dynamics --- Question 43.'' The supplied record lists Benedetto, DeMarco, Ingram, Jones et al. as proposers and gives the proposed year as 2018.
\item Supplied source attribution: Benedetto, DeMarco, Ingram, Jones et al., \emph{Current Trends and Open Problems in Arithmetic Dynamics} (2018), Question 43, \url{https://arxiv.org/abs/1806.04980}. The supplied background records access on 2026-10-16 and extraction by \texttt{arxiv\_latex}; these are supplied provenance claims, not independently verified access or extraction.
\item The local record labels the problem ``open'' and records an update on 2026-10-16. This is supplied registry status and dated metadata, not an independently verified assessment of the mathematical literature.
\end{itemize}
\end{document}
